\clearpage
\section{Технологическая часть}

\subsection{Средства реализации}

Алгоритмы реализованы на языке Свифт (Swift) версии 6.4.2. При выполнении
программы язык проверяет переполнение целочисленных операций и аварийно
завершает программу вместо получения неверного результата~\cite{swift}.
Для сборки используется диспетчер пакетов Свифт (Swift Package Manager),
для тестирования~--- библиотека Свифт Тестинг (Swift Testing)~\cite{swifttesting}.
Программа собиралась и выполнялась под управлением операционной системы
макОС (macOS) 26.6.2 на компьютере с процессором Apple M1 Pro и оперативной
памятью 16~ГБ.

\subsection{Реализация алгоритмов}

Реализации алгоритмов приведены в приложении~А. Нерекурсивный алгоритм
реализован функцией \texttt{countNumbersBySumDigits}
(листинг~\ref{lst:impl-iter} «Реализация нерекурсивного алгоритма»).
В ней цифры перебираются от 0 до 9 с досрочным выходом при отрицательной
сумме $j-c$, что равносильно границе $\min(9,j)$ в листинге~\ref{lst:iter}.

Функция \texttt{countNumbersBySumDigitsRecursive} реализует рекурсивный
алгоритм (листинг~\ref{lst:impl-rec} «Реализация рекурсивного алгоритма»).
Она вызывает вспомогательные функции,
соответствующие функциям $F$ и $S$ листинга~\ref{lst:rec}.

Для вывода таблицы значений программа использует функцию
\texttt{countNumbersBySumDigitsMatrix}, которая возвращает таблицу целиком
(листинг~\ref{lst:impl-table} «Формирование и вывод таблицы значений»).

\subsection{Сборка и запуск}

Команды сборки, запуска программы и тестирования выполняются из каталога
\texttt{lab-03} и приведены в листинге~\ref{lst:run} «Команды сборки, запуска
и тестирования».

\bigskip
\noindent\begin{minipage}{\dimexpr\linewidth-0.5mm\relax}
\begin{lstlisting}[caption={Команды сборки, запуска и тестирования},label={lst:run}]
swift build
swift run lab03
swift test
\end{lstlisting}
\end{minipage}
\medskip

При запуске программа выводит таблицу значений $F(i,j)$ для $s=15$, $k=5$
и результаты обоих алгоритмов: 2654.

\subsection{Тестирование}

Функциональные тесты приведены в листинге~\ref{lst:tests}
«Функциональные тесты» приложения~А. Одинаковые входные данные и ожидаемые
результаты для обеих реализаций приведены в таблице~\ref{tbl:tests}.

\begin{table}[H]
\caption{Функциональные тесты}
\label{tbl:tests}
\setlength{\extrarowheight}{2pt}
\begin{tabularx}{\linewidth}{|c|c|c|X|}
\hline
$s$ & $k$ & Ожидаемый результат & Проверяемый случай \\
\hhline{|=:=:=:=|}
4 & 1 & 1 & однозначное число \\ \hline
14 & 1 & 0 & сумма больше $9k$ для однозначных чисел \\ \hline
15 & 3 & 69 & пример из условия задачи \\ \hline
28 & 3 & 0 & сумма больше $9k$ \\ \hline
2 & 2 & 2 & числа 11 и 20 \\ \hline
1 & 4 & 1 & наименьшая сумма, число 1000 \\ \hline
36 & 4 & 1 & наибольшая сумма, число 9999 \\ \hline
\end{tabularx}
\end{table}

Все 14 проверок (по 7 для каждой реализации) пройдены успешно.
